\emph{Origin and numbering.} Clauses (i) to (iii) are clauses (i) to (iii) of Theorem synthesis.stationary-horizon-state-is-its-own-edge-scaling-limit, stated there for the generators of horizons in the scope of Premise (N$_{+}$); here they are stated on the carrier of one generator, and the arguments below are moved from that proof. The identification of each horizon's flow with the boost $\beta_{t}$, and the reading of (ii) as (N$_{+}$), stay there.

\emph{(i)} On fields the boost is $\beta_{t}(\psi(u)) = e^{\kappa t/2}\psi(e^{\kappa t}u)$, which is $\sigma^{0}_{\lambda}$ at $\lambda = e^{\kappa t}$. For $\sigma_{\lambda}f(u) = \lambda^{p}f(\lambda u)$ one has $\ip{\sigma_{\lambda}f}{\sigma_{\lambda}g} = \lambda^{2p - 1}\ip{f}{g}$, so the anticommutator $\{a(f), a^{*}(g)\} = \ip{f}{g}$ is preserved exactly for $p = \tfrac12$; the check verifies this on indicator functions of rational intervals at $\lambda = (3/2)^{4}, (1/5)^{4}$ among the weights $p = k/4$, $0 \le k \le 4$. In $\Cl(1,3)$ with $Q(e_{0}) = 1$, $Q(e_{j}) = -1$, $(e_{1}e_{0})^{2} = 1$ and $B = \cosh\tfrac{\lambda}{2} + \sinh\tfrac{\lambda}{2}\,e_{1}e_{0}$ is a rotor; $e_{1}e_{0}$ anticommutes with $e_{0}, e_{1}$ and commutes with $e_{2}, e_{3}$, and $e_{1}e_{0}\ell = \ell$, $e_{1}e_{0}n = -n$, so $B\ell\tilde B = B^{2}\ell = e^{\lambda}\ell$, $Bn\tilde B = e^{-\lambda}n$, and $u\ell + y \mapsto e^{\lambda}u\ell + y$ for $y$ in the span of $e_{2}, e_{3}$; the check verifies these at $e^{\lambda/2} = \tfrac32, \tfrac15, 7$. So the boost is the scaling of the generator about the edge point with the weight $\tfrac12$ of a chiral field of scaling dimension $\tfrac12$.

\emph{(ii)} If $\omega\circ\beta_{t} = \omega$ for all $t$, then by (i) $\omega\circ\sigma^{0}_{\lambda} = \omega$ for all $\lambda > 0$, that is $\lambda W_{\omega}(\lambda x, \lambda y) = W_{\omega}(x, y)$; the family whose limit defines the scaling limit at the edge is constant, so the limit exists and is $\omega$. If $\omega$ is locally the null vacuum at the edge, then $W_{\omega} = W_{0}$, and a gauge-invariant quasi-free state is determined by its two-point function, so $\omega = \omega_{0}$. Conversely $W_{0}$ is invariant under translations and under $\sigma^{u_{0}}_{\lambda}$ for every $u_{0}$, so $\omega_{0}$ is its own scaling limit at every point. By Theorem \claimref{synthesis.null-vacuum-on-a-horizon-generator}(ii) the gauge-invariant quasi-free ground states of the null translations are exactly $\omega_{0}$, which gives the stated equivalence. For the weaker form: $W_{0}$ is stationary, so $D = W_{\omega} - W_{0}$ satisfies $\lambda D(\lambda x, \lambda y) = D(x, y)$. If $|D| \le C$ on a neighbourhood $U$ of $(0, 0)$, then for any $(x, y)$ and $\lambda$ small enough that $(\lambda x, \lambda y) \in U$, $|D(x, y)| = \lambda|D(\lambda x, \lambda y)| \le \lambda C$, and $\lambda \to 0$ gives $D(x, y) = 0$.

\emph{(iii)} The two halves $\CAR(L^{2}(\R_{\pm}, du)\otimes K)$ generate $\CAR(L^{2}(\R, du)\otimes K)$ as a graded tensor product, and the product of two even states is an even state; each half is $\beta$-invariant, so the product of stationary states is stationary. On $u > 0$ put $u = e^{x}$ and $\phi(x) = e^{x/2}\psi(e^{x})$: the map is unitary from $L^{2}(\R_{+}, du)$ to $L^{2}(\R, dx)$, the boost becomes the translation $x \mapsto x + \kappa t$, and the null vacuum's restriction becomes $\tfrac{I_{4}}{2\pi}e^{(x + x')/2}(e^{x} - e^{x'} + I_{4}0)^{-1} = \tfrac{I_{4}}{2\pi}\big(2\sinh((x - x' + I_{4}0)/2)\big)^{-1}$, the chiral thermal two-point function at period $2\pi$ (Theorem \claimref{synthesis.null-vacuum-on-a-horizon-generator}(iii)). The chiral two-point function KMS at period $2\pi c$ for the translations is the same function of $(x - x')/c$ divided by $c$ (the leading singularity $\tfrac{I_{4}}{2\pi}(x - x' + I_{4}0)^{-1}$, fixed by the anticommutator, is then the same for every $c$), which gives $W_{c}$; its Fourier transform in $x - x'$ is that at $c = 1$ evaluated at $c\,w$, so the Mellin symbol is $m_{1}(c w) = (1 + e^{-2\pi c w})^{-1}$, the symbol of the I-state of Theorem \claimref{hawking.i-state-continuum} for $h = -\kappa w$ at $\beta = 2\pi c/\kappa$, which is jointly KMS by Fact \claimref{hawking.imported-quasi-free-kms-reduces-two}. In the check the kernel is evaluated exactly in the form $(2\pi/I_{4})W_{c} = (uv)^{-1/2}/\big(c(t^{1/2c} - t^{-1/2c})\big)$, $t = u/v$, at points $u = a^{2k}$, $v = b^{2k}$ with $k/c$ an integer ($c = \tfrac13, \tfrac12, 1, 2, 3$): $c = 1$ gives $1/(u - v)$; every kernel and every difference is homogeneous of degree $-1$ under scaling by squares; $D_{c} = 0$ exactly at $c = 1$; and along a ray towards the edge $D_{c}$ grows as the inverse scale; the symbol satisfies the two-point KMS identity at period $2\pi c$ and, for $c \neq 1$, not at $2\pi$; and $t^{1/2c} \mapsto -t^{1/2c}$ under $y \mapsto y - 2\pi I_{4}c$ makes the kernel antiperiodic there (the fermionic KMS condition), while at $c = \tfrac12$ the shift by $2\pi I_{4}$ leaves it periodic. For $c = \infty$ the symbol is $\chi_{(0, \infty)}(w)$, that is $\chi_{(-\infty, 0)}(h)$ and $(2\pi/I_{4})W_{\infty} = (uv)^{-1/2}/(y + I_{4}0)$.

With $S(z) = \sinh z/z$, $g_{c}(y) = y^{-1}\big(S(y/2c)^{-1} - S(y/2)^{-1}\big)$. Both inverses are even power series with constant term $1$, so the pole cancels; $S(z)^{-1} = 1 - \tfrac{z^{2}}{6} + \tfrac{7z^{4}}{360} - \cdots$ gives $g_{c}(y) = -\tfrac{1}{24}(c^{-2} - 1)y + \tfrac{7}{5760}(c^{-4} - 1)y^{3} + \cdots$, and the zeros of $\sinh(y/2c)$ and $\sinh(y/2)$ off $0$ lie at $2\pi I_{4}c k$ and $2\pi I_{4}k$, $k \neq 0$, so $g_{c}$ is analytic on $|\mathrm{Im}\,y| < 2\pi\min(c, 1)$. The check computes both series exactly to order $12$ in $y$ with $\gamma = c^{-2} \in \{0, \tfrac19, \tfrac14, 1, 4\}$: zero constant term, even, first coefficient $-(\gamma - 1)/24$, identically zero only at $\gamma = 1$; and as a control, omitting the factor $1/c$ leaves the pole $(c - 1)/y$. At a point $u_{0} > 0$, $\lambda W_{0}(u_{0} + \lambda x, u_{0} + \lambda y) = W_{0}(x, y)$ and $\lambda D_{c}(u_{0} + \lambda x, u_{0} + \lambda y) = O(\lambda)$ uniformly for $x, y$ in a compact set, because $D_{c}$ is smooth near $(u_{0}, u_{0})$; for small $\lambda$ both arguments stay on the half $u > 0$, so the cross-edge block does not enter. The scaling limit at $u_{0}$ is $\omega_{0}$, and the same holds at $u_{0} < 0$ by the reflection $u \mapsto -u$. Only the same-half block entered: for test functions supported in a compact set and $\lambda$ small, $u_{0} + \lambda x$ and $u_{0} + \lambda y$ lie on the half containing $u_{0}$. So the conclusion holds for every stationary gauge-invariant quasi-free state with the half restrictions of $\omega_{c}^{\times}$ and an arbitrary cross-edge block (subject only to positivity of the whole symbol), and for $c \neq 1$ such a state is not $\omega_{0}$, because its restriction to $u > 0$ is the I-state at period $2\pi c$ while that of $\omega_{0}$ is the I-state at $2\pi$ (Theorem \claimref{synthesis.null-vacuum-on-a-horizon-generator}(iii)) and the I-state at a given period is unique (Theorem \claimref{hawking.i-state-continuum}). Across the edge, $\omega_{c}^{\times}(\psi^{*}(v)\psi(u)) = 0$ for $u > 0 > v$ (product of even states, odd factors), while $W_{0}(u, v) = \tfrac{I_{4}}{2\pi}(u - v)^{-1} \neq 0$ there, homogeneous of degree $-1$: the difference is unbounded on every neighbourhood of $(0, 0)$, and the check tabulates it at three points and three scales. So $\omega_{c}^{\times} \neq \omega_{0}$ for every $c$; at $c = 1$ both halves are the I-state at period $2\pi$ (the continuum restatement of (R) on each half) and the state is still not the null vacuum.
